% Part: lambda-calculus
% Chapter: introduction
% Section: fixed-point-combinator

\documentclass[../../../include/open-logic-section]{subfiles}

\begin{document}

\olfileid{lam}{int}{fix}
\olsection{مؤلفات النقطة الثابتة}

لتكن لدينا عبارة لامبدا~$\OLId{latin-ordinary}{g}$، ولنطلب حدًا آخر~$\OLId{latin-ordinary}{k}$ يكون فيه~$\OLId{latin-ordinary}{k}$
مكافئًا بحسب~$\beta$ للحد~$gk$. نعرّف الحدين
\[
\fn{diag}(\OLId{latin-ordinary}{x}) = xx
\]
و
\[
\OLId{latin-ordinary}{l}(\OLId{latin-ordinary}{x}) = \OLId{latin-ordinary}{g}(\fn{diag}(\OLId{latin-ordinary}{x}))
\]
على مقتضى اصطلاحاتنا؛ فـ$\OLId{latin-ordinary}{l}$ هو الحد~$\lambd[\OLId{latin-ordinary}{x}][\OLId{latin-ordinary}{g}(xx)]$.
ثم نأخذ~$\OLId{latin-ordinary}{k}$ مساويًا للحد~$ll$، فنجد
\begin{align*}
\OLId{latin-ordinary}{k} & = (\lambd[\OLId{latin-ordinary}{x}][\OLId{latin-ordinary}{g}(xx)])(\lambd[\OLId{latin-ordinary}{x}][\OLId{latin-ordinary}{g}(xx)]) \\
& \red  \OLId{latin-ordinary}{g}((\lambd[\OLId{latin-ordinary}{x}][\OLId{latin-ordinary}{g}(xx)])(\lambd[\OLId{latin-ordinary}{x}][\OLId{latin-ordinary}{g}(xx)])) \\
& = gk.
\end{align*}
وعليه، إذا وضعنا
\[
\OLId{latin-looped}{Y} = \lambd[\OLId{latin-ordinary}{g}][((\lambd[\OLId{latin-ordinary}{x}][\OLId{latin-ordinary}{g}(xx)])(\lambd[\OLId{latin-ordinary}{x}][\OLId{latin-ordinary}{g}(xx)]))]
\]
كان~$Yg$ و$\OLId{latin-ordinary}{g}(Yg)$ قابلين للاختزال إلى حد واحد، فينتج
$Yg \equiv_\beta \OLId{latin-ordinary}{g}(Yg)$. وهذا هو «مؤلف كاري». وأما إذا وضعنا
\[
\OLId{latin-looped}{Y} = (\lambd[xg][\OLId{latin-ordinary}{g}(xxg)])(\lambd[xg][\OLId{latin-ordinary}{g}(xxg)])
\]
فإن~$Yg$ يختزل إلى~$\OLId{latin-ordinary}{g}(Yg)$ نفسه؛ وهذه نتيجة أقوى من مجرد التكافؤ.
وتسمى هذه الصيغة الأخيرة من~$\OLId{latin-looped}{Y}$ «مؤلف تورنغ».

\end{document}

